% Extracto íntegro por residencia del cierre físico-matemático sucesor.
\section{Respuesta autoinercial y principio de Mach desde Ambivalencia}
\label{sec:hmtf2-mach-ambivalencia}

Sea \(b:\mathcal X_\infty^{\rm enr}\to B_\infty\) la traza global de
frontera. Sobre \(B_\infty\) se conservan los lectores positivos
\(A(\beta)\), \(V(\beta)\), el reloj global \(Q(\beta)>0\) y la memoria
\(T(\beta)\). La bisagra de Ambivalencia proporciona
\[
 r_{\rm av}=\frac{\pi_{\rm HMT}}{\varphi_{\rm HMT}^2},
 \qquad
 r_{\rm av}^J=\frac{\varphi_{\rm HMT}}{\pi_{\rm HMT}^2}.
\]
Definimos pesos, sin datos gravitatorios externos,
\[
 w_A(\beta)
 =\frac{A(\beta)r_{\rm av}}
 {A(\beta)r_{\rm av}+V(\beta)r_{\rm av}^J},
 \qquad
 w_V(\beta)=1-w_A(\beta).
\]
Sean \(P_A,P_V\) proyecciones ortogonales complementarias del espacio local
de respuesta. Para energías internas \(E_1,E_2\), la respuesta
autoinercial HMT es
\[
 \overline I_v^{\rm HMT}(\beta)
 =\frac{E_1E_2}{Q(\beta)^2}
 \bigl(w_A(\beta)P_A+w_V(\beta)P_V\bigr),
\]
y
\[
 I_v^{\rm HMT}
 =\overline I_v^{\rm HMT}\circ b.
\]

\begin{theorem}[Cierre machiano de la respuesta de Ambivalencia]
\label{thm:hmtf2-respuesta-mach}
La respuesta \(I_v^{\rm HMT}\) es autoadjunta, positiva y constante en cada
fibra de \(b\). Por tanto, factoriza de modo único por la totalidad global.
La involución de Ambivalencia que intercambia
\[
 (A,r_{\rm av},P_A)\longleftrightarrow
 (V,r_{\rm av}^J,P_V)
\]
conjuga \(\overline I_v^{\rm HMT}\) sin cambiar su espectro. Si existen
\(x,x'\) con la misma restricción local, \(b(x)\ne b(x')\) y distinto
\(Q,A\) o \(V\), entonces la respuesta no factoriza sólo por la restricción
local y es machiana en sentido fuerte.
\end{theorem}

\begin{proof}
Los pesos son positivos y suman uno; la combinación de proyecciones es
positiva y autoadjunta. La fórmula sólo depende de \(b(x)\), luego es
constante en sus fibras y aplica el criterio de factorización por la
totalidad. El intercambio permuta los dos bloques espectrales. La última
afirmación es precisamente el control negativo frente a una respuesta
puramente local.
\end{proof}


La construcción posee además una lectura geométrica recuperada. Para el
estado conjunto \(\Gamma=(x,\mathcal O,m)\), sea
\(\nabla^{\rm tot}_{b(x)}\) la conexión inducida por la frontera y sea
\(\pi_S\) la proyección al subsistema. Si
\[
 \nabla^{\rm tot}_{\dot\Gamma}\dot\Gamma=0,
\]
la aceleración publicada es
\[
 a_S
 =\nabla^S_{\dot\gamma_S}\dot\gamma_S
 =\mathcal K_{\pi_S}^{\,b(x)}(\dot\Gamma,\dot\Gamma).
\]
Como el defecto depende de la conexión inducida por \(b(x)\), también
factoriza por la totalidad. Esto cierra matemáticamente el reemplazo de la
dicotomía inercial/no inercial por sistemas autoinerciales: el estado total
es geodésico y la aceleración local es defecto de proyección. La igualdad de
esta respuesta con una aceleración gravitatoria observada pertenece a
\textsc{reconocimiento externo}.

